\pdfvariable trailerid {[<C3192026090300000000000000000000><C3192026090300000000000000000000>]}
\documentclass[10pt]{article}
\usepackage[margin=0.72in]{geometry}
\usepackage{amsmath,amssymb,amsthm,booktabs,microtype,hyperref}
\hypersetup{colorlinks=true,linkcolor=blue,urlcolor=blue}
\providecommand{\CRevisionRound}{2}
\newtheorem{theorem}{Theorem}
\newtheorem{corollary}[theorem]{Corollary}
\newcommand{\RR}{\mathbb R}
\setlength{\emergencystretch}{3em}
\title{Ancient Clifford-Product Mean-Curvature Flow:\\Exact Lifespans, Focal Cylinders, and Minimal Stability}
\author{Route-A source-local certificate HCS-C319}
\date{3 September 2026}

\begin{document}
\maketitle
\begin{abstract}
For every $p,q\geq1$, we solve mean-curvature flow inside the unit sphere
for the full parallel family $S^p(\cos\theta)\times S^q(\sin\theta)$.
The affine coordinate $y=\sin^2\theta$ gives both ancient branches, their
exact finite lifespans, and their focal collapses.
\ifnum\CRevisionRound>0
We identify both Type-I round-cylinder blow-ups and prove exact area
dissipation.
\fi
\ifnum\CRevisionRound>1
At the minimal leaf we compute the entire separated Jacobi spectrum below
and at zero, obtaining index $p+q+3$ and nullity $(p+1)(q+1)$.
\fi
\end{abstract}

\section{Exact phase portrait and lifespan}
Put $n=p+q$, and write
\[
 \Sigma_\theta=S^p(\cos\theta)\times S^q(\sin\theta)
 \subset S^{n+1},\qquad 0<\theta<\frac\pi2.
\]
For $x\in S^p$, $z\in S^q$, choose
$\nu=(-\sin\theta\,x,\cos\theta\,z)$.  Our scalar convention is
$\partial_tF=H\nu$; this is the negative area gradient for the signs below.

\begin{theorem}[Complete product-family flow]\label{thm:phase}
The principal curvatures are $\tan\theta$ with multiplicity $p$ and
$-\cot\theta$ with multiplicity $q$.  Hence
\[
 H=p\tan\theta-q\cot\theta,
 \qquad y'=2(ny-q),\qquad y=\sin^2\theta.
\]
The unique stationary leaf has $y_*=q/n$.  Every other solution is ancient,
\begin{equation}\label{eq:solution}
 y(t)=\frac qn+\left(y_0-\frac qn\right)e^{2nt},
\end{equation}
and has one finite forward endpoint.  If $y_0<q/n$, then
\[
 T_-=\frac1{2n}\log \frac q{q-ny_0},\qquad y(t)\downarrow0,
\]
and the $S^p$ factor survives as the focal submanifold.  If $y_0>q/n$, then
\[
 T_+=\frac1{2n}\log \frac p{ny_0-q},\qquad y(t)\uparrow1,
\]
and the $S^q$ factor survives.  In both cases $y(t)\to q/n$ as
$t\to-\infty$.
\end{theorem}
\begin{proof}
Differentiating $F=(\cos\theta\,x,\sin\theta\,z)$ gives
$\partial_\theta F=\nu$.  The two shape eigenvalues follow by differentiating
$\nu$ along each factor.  Thus $\theta'=H$, and multiplication by
$2\sin\theta\cos\theta$ yields the affine equation for $y$.  Solving it
gives~\eqref{eq:solution}; imposing $y=0$ or $y=1$ gives the stated times.
Their logarithm arguments exceed one on the corresponding branches.
Backward convergence, uniqueness of the stationary leaf, and the focal
labels now follow directly.
\end{proof}

The faces $y=0,1$ are focal submanifolds, not regular product
hypersurfaces.  The cases $p=0$ or $q=0$ are excluded degenerate
sphere/double-cover geometries.  The backward endpoint is at infinite time,
not an added slice.

\ifnum\CRevisionRound>0
\section{Type-I focal-cylinder geometry}
The exact curvature and area density (up to the constant sphere volumes) are
\begin{align}
 |A|^2&=p\frac y{1-y}+q\frac{1-y}y,\label{eq:A2}\\
 \mathcal A(y)&=(1-y)^{p/2}y^{q/2}.\label{eq:area}
\end{align}

\begin{theorem}[Two Type-I limits and strict dissipation]\label{thm:typeI}
Every nonstationary solution of Theorem~\ref{thm:phase} satisfies
\[
 \frac d{dt}\log\mathcal A=-H^2<0.
\]
At a left collapse, $y=2q(T_--t)+O((T_--t)^2)$ and
\[
 (T_--t)|A|^2\longrightarrow\frac12,\qquad
 \text{blow-up cylinder }S^q(\sqrt{2q})\times\RR^p.
\]
At a right collapse,
$1-y=2p(T_+-t)+O((T_+-t)^2)$, the same residue is $1/2$, and the cylinder is
$S^p(\sqrt{2p})\times\RR^q$.
\end{theorem}
\begin{proof}
Taking the logarithmic derivative of~\eqref{eq:area} and substituting $y'$
gives
\[
 \left(\frac q{2y}-\frac p{2(1-y)}\right)2(ny-q)
 =-\frac{(ny-q)^2}{y(1-y)}=-H^2.
\]
Linearising the exact affine solution at either finite endpoint gives the
two expansions.  Equation~\eqref{eq:A2} then gives the Type-I residues.
Under parabolic rescaling by $(T-t)^{-1/2}$, the collapsing sphere has
squared radius $2q$ on the left and $2p$ on the right; the noncollapsing
factor becomes its Euclidean tangent space.  This proves the ordered
cylinder statements.
\end{proof}
\fi

\ifnum\CRevisionRound>1
\section{Jacobi index, nullity, and Route-A boundary}
At the stationary leaf the radii are $\sqrt{p/n}$ and $\sqrt{q/n}$,
$|A|^2=n$, and the normal Jacobi operator is
$J=\Delta+|A|^2+n=\Delta+2n$.  We fix the second-variation convention
$\mathcal Q(f,f)=-\int_{\Sigma}fJf$.

\begin{theorem}[Complete minimal stability count]\label{thm:spectrum}
For spherical-harmonic degrees $\ell,m\geq0$, the eigenvalue of $-\Delta$ is
\[
 \lambda_{\ell m}=\frac{n\ell(\ell+p-1)}p
 +\frac{nm(m+q-1)}q.
\]
Exactly the sectors $(0,0),(1,0),(0,1)$ lie below $2n$, and only $(1,1)$
equals $2n$.  Therefore, for the quadratic-form convention associated to
$J$, the Morse index is $n+3$ and the nullity is $(p+1)(q+1)$.
\end{theorem}
\begin{proof}
Separation of variables on the two round factors gives $\lambda_{\ell m}$.
For one factor, $\ell(\ell+p-1)/p$ equals $0$ at $\ell=0$, equals $1$ at
$\ell=1$, and exceeds $2$ for $\ell\geq2$; the same holds for $m,q$.
Thus the stated four sectors exhaust the values at or below $2n$.
Their harmonic multiplicities are $1,p+1,q+1$, and
$(p+1)(q+1)$, proving the counts.
\end{proof}

\paragraph{Finite evidence.}
The content-addressed atlas records 100 dimension pairs, 600 branch rows,
and 3,600 spectral cells: 25,858 scalar leaves.  An implementation-independent
checker performs 22,738 checks, including independent normalized-area and
dissipation calculations; SymPy closes 1,204 identities.  Isolated replay
is byte exact and 39 repaired-hash/parser mutations are rejected.  These
grids are regressions; Theorems~\ref{thm:phase}--\ref{thm:spectrum} are
analytic for all declared parameters.

\paragraph{Collision boundary.}
C281 treats intrinsic homogeneous Ricci flow, while the present evolution
is extrinsic spherical mean-curvature flow.  C314 concerns planar curve
shortening, whereas the present theorem closes a higher-dimensional
isoparametric hypersurface family and its minimal Jacobi spectrum.

\paragraph{Route-A result and nonclaims.}
Under \texttt{NO\_BAD\_EULER\_OR\_ROOT\_NUMBER}, the strict tuple is
\[
 (\mathrm{A0\_FAIL},\mathrm{A1\_FAIL},\mathrm{A2\_FAIL},
  \mathrm{A3\_FAIL},\mathrm{A4\_FORMAL\_HINT}).
\]
The model supplies no arithmetic-prime owner, primitive-orbit Euler ledger,
Euler product, target functional equation, counting law, or target-zero
match.  The Jacobi operator is source-local and only a formal spectral hint;
it is not a Hilbert--P\'olya operator.  Route A is rejected and Route B stays
locked.  No classification beyond the two-factor family, target local data,
root number, or automorphy is claimed.

\paragraph{AI use.}
A generative language model assisted drafting and code scaffolding.  The
displayed proofs, independent recomputation, hostile tests, and deterministic
artifacts define the audit record.
\fi

\section*{Source lineage}
\begin{thebibliography}{9}
\bibitem{ReisTenenblat}
F. dos Reis and K. Tenenblat,
``The mean curvature flow by parallel hypersurfaces,''
\emph{Proceedings of the American Mathematical Society} 146 (2018),
4867--4878. DOI:
\href{https://doi.org/10.1090/proc/14178}{10.1090/proc/14178}.

\bibitem{LiuTerng}
H.-L. Liu and C.-L. Terng,
``The mean curvature flow for isoparametric submanifolds,''
\emph{Duke Mathematical Journal} 147 (2009), 157--179. DOI:
\href{https://doi.org/10.1215/00127094-2009-009}%
{10.1215/00127094-2009-009}.
\end{thebibliography}
\end{document}
